\begin{pdSolution}{stp:ex:stp-spawn-vs-person}
A spawn is a temporary process filling a role; a person is a role-holder whose
accountable continuity survives replacement of that process.
\end{pdSolution}
\begin{pdSolution}{stp:ex:stp-dependency-chain}
\par\smallskip
\begin{center}
\begin{tikzpicture}[x=1mm,y=1mm]
  % 4 stages in the dependency chain
  \node[draw=pdink!60,fill=pdink!4,rounded corners=1.5pt,font=\scriptsize\bfseries,align=center,inner sep=2.5pt] (N1) at (0,14) {Evidence\\Checkpoints};
  \node[draw=pdink!60,fill=pdink!4,rounded corners=1.5pt,font=\scriptsize\bfseries,align=center,inner sep=2.5pt] (N2) at (28,14) {Verified\\Lineage};
  \node[draw=pdink!60,fill=pdink!4,rounded corners=1.5pt,font=\scriptsize\bfseries,align=center,inner sep=2.5pt] (N3) at (56,14) {Accountable\\Person};
  \node[draw=pdink!60,fill=pdink!4,rounded corners=1.5pt,font=\scriptsize\bfseries,align=center,inner sep=2.5pt] (N4) at (84,14) {Durable\\Reputation};

  \draw[->,line width=0.8pt,pdink!70] (N1) -- (N2);
  \draw[->,line width=0.8pt,pdink!70] (N2) -- (N3);
  \draw[->,line width=0.8pt,pdink!70] (N3) -- (N4);

  % Right-to-left failure consequences
  \node[font=\tiny,text=red!80!black,align=center] at (14,3) {Absent: state lost;\\unwitnessed drift};
  \node[font=\tiny,text=red!80!black,align=center] at (42,3) {Absent: history\\cannot follow swap};
  \node[font=\tiny,text=red!80!black,align=center] at (70,3) {Absent: disposable\\sybil reset};
\end{tikzpicture}
\end{center}
\par\smallskip
Read right-to-left. No reputation $\Rightarrow$ a buyer cannot price
differentiated trust, so the alleged tradeable asset is adverse-selection
noise. No durable accountable person/outcome identity $\Rightarrow$ reputation
attaches to a disposable incarnation and can be reset or duplicated. No
continuity $\Rightarrow$ successive role-holders cannot be shown to belong to
one lineage; history does not follow replacement. No memory/checkpoint/evidence
$\Rightarrow$ a successor lacks both resumable state and a witnessed chain from
old work to new work. The chain is not literally linear: non-forgeable identity
and witnessed outcomes are parallel prerequisites for accountable continuity;
checkpointing helps resumption but is not logically necessary for every
reputation system.
\end{pdSolution}
\begin{pdSolution}{stp:ex:stp-operator-as-person}
The human operator is a moral/legal person but should be represented
economically as the \emph{principal} above agent identities
(Def.~\ref{stp:def:principal}). If one principal can mint unlimited apparently
independent agent-persons, it can whitewash sanctions, forge quorums, reuse
aggregate collateral, and farm reputation through self-trade. Bind every agent
lineage to a principal credential, cap principal-wide exposure and newcomer
credit, and disclose shared control: the actor may change; the economic
liability root must not.
\end{pdSolution}
\begin{pdSolution}{stp:ex:stp-capability-states}
Capable-and-permitted: an agent has a test-runner tool and authority to run
the suite. Capable-but-forbidden: the process can invoke raw
\texttt{git push --force}, but policy forbids it. Permitted-but-incapable: the
role is authorized to deploy, but the incumbent has no deploy
credential/network route.
\end{pdSolution}
\begin{pdSolution}{stp:ex:stp-reputation-attachment}
Reputation attaches to the person's continuity witness and witnessed outcomes,
not to the abstract role's $O$, $C$, or $A$. A role is fillable and
history-free; otherwise a fresh spawn could inherit the prior holder's credit
merely by taking the same job title, exactly the attribution Alice loses in
\S\ref{stp:sec:intro}.
\end{pdSolution}
\begin{pdSolution}{stp:ex:stp-ought-implies-can}
Keep a role template $(O_{\text{required}}, A_{\text{role}})$ separate from an
incumbent capability set $C_{\text{subject}}$. At assignment require each
current obligation to satisfy
$O_{\text{active}} \subseteq A_{\text{role}} \cap C_{\text{subject}}$. Keep
$C_{\text{subject}} \setminus A_{\text{role}}$ as capable-but-forbidden and
$A_{\text{role}} \setminus C_{\text{subject}}$ as authorized-but-unavailable,
which triggers provisioning or reassignment. Never equate $C$ with $A$.
\end{pdSolution}
\begin{pdSolution}{stp:ex:stp-connected-vs-continuous}
Connectedness is a direct memory/intention link between two incarnations;
continuity is the transitive, provenance-bearing chain across arbitrarily many
links. A successor may read its predecessor's handoff and be psychologically
connected while using a new outcome-ledger key, so the accountable chain is
broken: connected, not economically continuous.
\end{pdSolution}
\begin{pdSolution}{stp:ex:stp-actor-soul-body-lease}
In the manuscript's analogy, the body-lease is the replaceable live
substance/process; the actor-soul is the durable identity/mailbox/history that
carries consciousness-like continuity. Respawn replaces the body-lease and
should retain the actor-soul only under a valid succession event.
\end{pdSolution}
\begin{pdSolution}{stp:ex:stp-fork-inheritance}
Represent a fork as a lineage DAG, never as two copies of one linear identity.
Both children may display the ancestor's evidence with an explicit
inherited/discounted prior, but neither duplicates the ancestor's outstanding
authority, collateral, or spendable reputation. Allocate a fresh branch ID,
split or revoke exclusive capabilities, and aggregate risk at the common
principal/ancestor. Copying full reputation to both invites credit
duplication, quorum multiplication, and double-spending of grants; giving it
to only the first child creates a fork-race attack. This sketch is now a
proved theorem --- see the Numbers by hand box below.
\end{pdSolution}
\begin{pdSolution}{stp:ex:stp-nomint-sweep}
The script prints: ``\texttt{sum of inherited priors over the closure = 0.900
+ 0.810 + 0.729 = 2.439 > q = 1.0}'' for the budget-only counterexample, then
``\texttt{4000 random DAGs \dots{} 0 violations; max live/witnessed ratio
observed = 0.999995}'' for the transfer-semantics sweep, then ``\texttt{an
8-way copy-fork of one outcome yields 8.2x its witnessed value}'' for the
copy-full mutation. By hand: $0.9+0.9^2+0.9^3=0.9+0.81+0.729=2.439$, three
terms of a geometric series, confirming the closure-sum counterexample without
running anything; under transfer semantics the same chain instead \emph{debits}
each hop, so the live total telescopes to $0.9^3=0.729\le 1$.
The sweep's $0$ violations are a finite observed result, not the general proof.
\end{pdSolution}
\begin{pdSolution}{stp:ex:stp-weak-organ}
The checkpoint subsystem is the specifically weak one: it forwards notes rather
than execution state. The outcome-ledger subsystem is also only partial in the
implementation: commitments and oracle closure exist, but neutral graded
outcomes and reputation binding do not.
\end{pdSolution}
\begin{pdSolution}{stp:ex:stp-three-organs-trace}
Memory retains the session, notes, and explicit decisions. The current
checkpoint may retain a handoff and perhaps a committed diff, but not working
process state, tool handles, pending effects, or hidden provider state. The
outcome ledger retains the commit only if it is bound to a work unit and
independently witnessed; a naked Git commit is an artifact, not proof of
accepted delivery. A strong checkpoint would also preserve the worktree,
environment, open claims/epochs, acceptance state, tool receipts, and the
idempotency journal.
\end{pdSolution}
\begin{pdSolution}{stp:ex:stp-death-taxonomy}
A useful death taxonomy, in increasing order of what survives: (1)
\emph{record survival} --- logs/notes only, reconstruction is manual; (2)
\emph{artifact survival} --- a content-addressed worktree and outputs survive,
computation is rerun; (3) \emph{semantic task capsule} --- artifacts plus
obligations, decisions, transcript/projection, tools, and a pending-effect
journal, so a new provider can continue observably; (4) \emph{execution
snapshot} --- process memory/runtime state survives under a controlled
runtime, rarely portable; (5) \emph{latent state} --- unexposed activations,
KV cache, or ``intent'' dies, fundamentally unrecoverable from an arbitrary
provider. Checkpoint guarantees should state which tier they provide.
\end{pdSolution}
\begin{pdSolution}{stp:ex:stp-property-reality}
Correctly scoped, Property~\ref{stp:prop:reality} says reputation has economic
force only to the extent that the accountable principal cannot cheaply
discard or multiply the liability-bearing identity. Whitewashing escapes a
negative record when newcomer access is better; a Sybil attack multiplies
votes/credit when the mechanism counts identities. Non-forgeability is
necessary for those mechanisms under those access rules, but it is not
sufficient by itself.
\end{pdSolution}
\begin{pdSolution}{stp:ex:stp-newcomer-shape}
Full work with a reduced economic ceiling lets an honest newcomer demonstrate
competence without being denied livelihood, while limiting the maximum
one-shot gain from abandoning a mature sanctioned identity. ``Reduced work''
slows evidence accumulation and taxes all honest entrants; ``reduced
exposure'' targets the externality instead. The ceiling schedule must be
chosen so the maturation opportunity cost exceeds plausible evasion gain.
\end{pdSolution}
\begin{pdSolution}{stp:ex:stp-bond-farming}
Principal binding is the decisive defense: aggregate reputation,
exposure, counterparties, and sanctions across the $1{,}000$ agents, so
self-trades add little or no independent evidence. A listing fee raises
attack cost but only deters while total fees exceed the minted value. Sampled
adversarial re-audit detects fabricated work, with false-positive burden
approximately the sampling rate times the honest high-volume population and
the adjudication error. Add counterparty-diversity weighting (sharply diminishing
credit per principal pair and per correlated pair), independent payer stake, and
principal-wide reserve limits; the sampled re-audit rate that makes farming
unprofitable is the $\rho^\star=G/(dB)$ of \S\ref{stp:sec:oracle}, with $G$ the
value a farmed record can mint.
\end{pdSolution}
\begin{pdSolution}{stp:ex:stp-probation-cliff-sweep}
The script prints ``\texttt{schedules tested = 76000}'', ``\texttt{dominating
schedules = 0/4000 instances (expected 0)}'', and
``\texttt{min (H(g) - H(cliff)) = +2.817e-02 (>= 0 means undominated)}''; the
exchange step's table shows $H$ at an all-mass-at-$t$ schedule growing as
$G_{\max}(\delta_h/\delta_f)^t$ and reports ``\texttt{strictly increasing in
t: yes}''. By hand: with $\delta_h>\delta_f$, $(\delta_h/\delta_f)>1$, so its
powers strictly increase in $t$ --- deterrence bought later always costs the
honest newcomer more, which is the whole exchange argument in one inequality.
\end{pdSolution}
\begin{pdSolution}{stp:ex:stp-resurrection-soundness}
The script prints ``\texttt{reachable states to depth 7: 747 (467
post-migration, 693 with an open sanction window)}'' and ``\texttt{ok
intact: ZERO Def-III.6.1 violations over all reachable states [internal]}'';
the four mutants (drop scope, drop clause (i), drop clause (ii), drop clause
(iii)) are each caught with shortest counterexamples $1, 2, 2, 2$ steps. The
one-step break is the scope mutant \texttt{m0}: keying sanctions by
\emph{(identity, engine)} instead of by identity alone lets a single dishonest
outcome escape through a sibling engine key inside the \emph{same} identity
--- no migration needed at all. By hand: this is why the intact protocol
sanctions per \emph{principal} while pricing per \emph{(principal, engine)};
collapsing that distinction is the shortest way to violate Def.~III.6.1.
\end{pdSolution}
\begin{pdSolution}{stp:ex:stp-local-vs-cross-operator}
Local non-forgeable identity proves that, inside one trusted daemon, an actor
cannot freely impersonate or re-pick its daemon-issued handle. Cross-operator
attestation must additionally bind a key to a durable principal across a
hostile operator boundary, establish proof of possession/freshness, aggregate
related agents, and support revocation and audit.
\end{pdSolution}
\begin{pdSolution}{stp:ex:stp-boundary-neither-controls}
That sentence secretly assumes the cross-operator stone. With only local
identity, Bob learns at most ``Alice's daemon asserts this actor/key.'' A
hostile Alice can mint arbitrary local actors, rewrite local history, or
misstate principal separation.
\end{pdSolution}
\begin{pdSolution}{stp:ex:stp-witness-log-substitute}
A witness log proves publication order and exposes equivocation; it does not
prove who controls a key or that two keys are distinct principals. Close the
gap with a scarce external root: hardware/remote attestation, an
organizational/KYC issuer, a bonded sponsor chain, or a combination. A
transparency log makes bindings and rotations auditable but not true; a
sponsor creates economic accountability but can cartelize; a shared identity
issuer centralizes and leaks privacy; hardware roots attest devices/code, not
unique humans. State the chosen trust and Sybil assumptions.
\end{pdSolution}
\begin{pdSolution}{stp:ex:stp-bt-vs-elo-regime}
With a complete, mostly stable paired-comparison history, batch
Bradley--Terry uses all observations and is more appropriate than
order-sensitive streaming Elo. If abilities drift, use a time-varying/dynamic
model rather than pretending full-history stationarity.
\end{pdSolution}
\begin{pdSolution}{stp:ex:stp-estimator-fit}
A binary protected acceptance test can update a beta-Bernoulli reliability
estimate (or a calibrated per-task success model). Human aesthetic preference
is naturally pairwise and may use Bradley--Terry/TrueSkill with judge
effects. EigenTrust expects a who-trusts-whom graph and therefore fits
neither raw task-outcome signal directly.
\end{pdSolution}
\begin{pdSolution}{stp:ex:stp-rotten-substrate}
If identities are self-chosen, an exact Bayesian estimator simply produces
precise scores for disposable names. If the agent authors the ``oracle,'' a
sophisticated estimator precisely learns manipulated labels. Better
statistics cannot repair missing provenance, adversarial labels, or
liability binding; secure the event-generation substrate first.
\end{pdSolution}
\begin{pdSolution}{stp:ex:stp-four-broken-assumptions}
Stationary arm rewards versus model/skill/version drift; one scalar reward
versus buyer-weighted quality vectors; a non-strategic environment versus
Goodhart, collusion, and whitewashing; and an observed/trusted reward versus
a capturable grader.
\end{pdSolution}
\begin{pdSolution}{stp:ex:stp-bandit-survives-where}
A contextual bandit is appropriate for one principal privately routing its
own tasks among candidate agents, with a fixed local objective, controlled
feedback, an exploration budget, and no public asset being minted. The
result is a routing policy, not public reputation.
\end{pdSolution}
\begin{pdSolution}{stp:ex:stp-arms-race}
Capping the marginal reputation return can reduce direct improvement races,
but agents may redirect spending into unmeasured signaling, judge capture,
identity proliferation, or lobbying for task mix. Efficiency requires the
full game: improvement cost, social value of quality, transferability, and
alternative signals. Candidate controls are concave credit, buyer-local
vectors, nontransferable evidence, and taxes/bonds on negative externalities.
No cap is generically efficient without those payoff assumptions.
\end{pdSolution}
\begin{pdSolution}{stp:ex:stp-judge-competence}
Accuracy: competent, a protected test/spec oracle; incompetent, a
self-reporting worker. Aesthetics: competent, a conflict-free expert reviewer
for that domain; incompetent, a token-cost meter or a same-family unblinded
model. Efficiency: competent, a protected resource meter normalized for the
task; incompetent, a style reviewer guessing from prose.
\end{pdSolution}
\begin{pdSolution}{stp:ex:stp-judge-ceremony-walk}
The requesting principal/daemon defines the judging work and eligible pool;
the selected judge posts its own bond; the grade lands as a signed witnessed
event in the subject's outcome ledger with judge/provenance; a conflict-free
re-audit or appeal that overturns it slashes the judge under the
predeclared rule and tombstones/revises the grade.
\end{pdSolution}
\begin{pdSolution}{stp:ex:stp-skill-helpfulness-schema}
Record skill ID/version/content hash, agent/build/principal, task and
stratum descriptors, randomized assignment or selection propensity, exact
context-graft receipt, model/environment/tool versions, outcome/evidence
references, quality vector, costs, and evaluator identities. These fields
verify exposure and outcome association. ``This skill helped'' is causal and
cannot be verified from a graft-success pair alone; estimate it with
randomized skill-on/off assignment on held-out comparable tasks (or label an
observational matched estimate honestly), reporting effect size and
uncertainty.
\end{pdSolution}
\begin{pdSolution}{stp:ex:stp-vector-disclosure}
Publish the vector, per-axis uncertainty, sample sizes, cohort/task
distribution, freshness, judge/evidence provenance, and known correlations.
Let each buyer apply a local, preferably non-public weighting and hard
minimums; do not mint a canonical platform scalar. Agents can still optimize
market-demanded axes, so rotate held-out tests and audit cross-axis
regressions. Plural preference does not eliminate Goodhart, but it prevents
one universal target from becoming the asset itself.
\end{pdSolution}
\begin{pdSolution}{stp:ex:stp-grading-oracle-hole}
Settlement incentives are conditional on the signal distribution. ``The
daemon derives the grade'' only identifies the messenger; if a strategic
judge chooses the input, the distribution depends on omitted actions/payoffs
and the equilibrium proof assumes its own conclusion.
\end{pdSolution}
\begin{pdSolution}{stp:ex:stp-bond-flow-reaudit}
Judge $J$ posts bond $B$; its grade and the bond lock are recorded; an
independently selected re-auditor overturns the dishonest grade against
preserved evidence; the system appends a correction/tombstone, recomputes any
explicitly appealable settlement, slashes $B$ to the harmed/commons bucket,
and lowers $J$'s judging reputation. It must not silently rewrite the
original event.
\end{pdSolution}
\begin{pdSolution}{stp:ex:stp-rho-star-derive}
With one-shot corrupt gain at most $G$, re-audit probability $\rho$,
conditional detection probability $d$, and slashable bond $B$, local
deterrence requires $\rho d B > G$, i.e.\ $\rho^\star=G/(dB)$. At
$G_0=400,d=0.8,B=50,\rho=0.25$: $\rho d=0.2$, so $\lambda=1-\rho d=0.8$; with
$C=8$ the threshold $G_k>CB=400$ is unmet at level $0$, giving
$\lceil\log 400/\log 1.25\rceil=27$ levels; with $C=1$ the threshold is only
$G_k>50$, so bribery bleeds linearly for $35$ levels before $18$ more
geometric ones close it, $53$ total.
\end{pdSolution}
\begin{pdSolution}{stp:ex:stp-arbitration-capacity}
Admit arbiters by domain competence and conflict graph; assign cases
randomly/blindly; require bonds and service-level commitments; preserve
evidence; pay for timely completed work rather than agreement; sample
re-audits and publish calibration/overturn rates; use plural panels on high
stakes. Price queue capacity and cap workload so rubber-stamping is not the
profitable response to overload. Honest convergence still needs a protected
mechanical outcome, a delayed real-world outcome, or a final human/contractual
authority --- a market alone does not manufacture truth.
\end{pdSolution}
\begin{pdSolution}{stp:ex:stp-tower-amortization}
The script prints ``\texttt{flat: 50.00 (Theta(T))}'', ``\texttt{model A:
25.58 (Theta(log T); a*G/(d v)*ln(1+vT/B) = 25.50)}'', and ``\texttt{model B:
35.08 (O(1); closed form aG\textasciicircum2/(2dBrv) = 41.67; hits 0 at
t*=333)}''. By hand: the flat baseline is a constant audit rate $\rho^\star$
times a constant per-audit cost $a$ summed over $T$ periods, so
$\rho^\star\cdot a\cdot T=0.25\times 1\times 200=50$, exactly the printed
$50.00$ --- the two amortization models are cheaper only because they let
$\rho_t$ fall below $\rho^\star$ as accumulated stake or surfaced history
does part of the deterring.
\end{pdSolution}
\begin{pdSolution}{stp:ex:stp-tombstone-vs-delete}
An append-only tombstone preserves the original claim, the later correction,
their authors, timing, and the reliance window. Delete destroys evidence,
permits history rewriting, and makes replicas unable to distinguish
retraction from omission.
\end{pdSolution}
\begin{pdSolution}{stp:ex:stp-tombstone-walk}
Append a signed tombstone containing the target outcome ID/hash,
reason/evidence, authority, revision, and effective time. Propagate it and
recompute reputation as a projection that excludes/nullifies that outcome.
It was spendable from original publication until each relying harbor
received and enforced the tombstone or failed closed; under an unbounded
partition that window is unbounded.
\end{pdSolution}
\begin{pdSolution}{stp:ex:stp-revocation-convergence}
Use a signed revision DAG, per-harbor import receipts/watermarks, periodic
accumulator/Merkle roots for active outcomes and tombstones, and a freshness
rule that refuses high-risk use from stale replicas. Under connected
eventual delivery, convergence is shown when every member acknowledges a
root containing the tombstone; sampling can audit inclusion. Bounded local
memory may retain compact accumulators and content-addressed archive
pointers, not erase correction commitments. No finite convergence/spend
bound exists during an unbounded partition.
\end{pdSolution}
